\subsection{复数的无穷和}

由于域 C 的加法群与加法群 \(R^{2}\) 相同，无须对 C 中的可和族与级数作专门研究，因为这已包含在第七章 § 3 的一般理论之中；我们把这理论的结果翻译成复数理论的语言这一练习留给读者。我们只叙述下面的命题，它是第七章 § 3 第 1 号命题 3 的一个推论：

命题 1 。若 \((u_{\lambda})_{\lambda \in \mathbf{L}}\) 与 \((v_{\mu})_{\mu \in \mathbf{M}}\) 是两个复数的可和族，则族 \((u_{\lambda}v_{\mu})_{(\lambda ,\mu)\in \mathbf{L}\times \mathbf{M}}\) 可和，且有

\[
\sum_ {(\lambda , \mu) \in \mathbf {L} \times \mathbf {M}} u _ {\lambda} v _ {\mu} = \left(\sum_ {\lambda \in \mathbf {L}} u _ {\lambda}\right) \left(\sum_ {\mu \in \mathbf {M}} v _ {\mu}\right).\tag{1}
\]

把关于四元数的相应结果的叙述留给读者。

